\documentclass[11pt]{article}
\usepackage[margin=0.8in]{geometry}
\usepackage[T1]{fontenc}
\usepackage{lmodern}
\usepackage{pdflscape}
\usepackage{longtable}
\usepackage{booktabs}
\usepackage{array}
\usepackage{threeparttable}
\usepackage{caption}
\usepackage{microtype}
\usepackage{amsmath,amssymb}
\usepackage[hidelinks]{hyperref}

\setlength{\parindent}{0pt}
\setlength{\parskip}{0.6em}
\renewcommand{\thetable}{S\arabic{table}}

\begin{document}

\begin{center}
{\Large\bfseries Supplementary Material}\\[0.5em]
{\large Sex Differences in the Association Between Insulin Resistance and Systemic Inflammation}\\
{\large Below the HbA1c Prediabetes Threshold}\\[0.5em]
Scott Moerschbacher; Janice Baker
\end{center}

\setcounter{table}{0}

% Supplementary statistical tables for Study #2.
% These tables use frozen analysis results and do not introduce new models.

\begin{table}[htbp]
\centering
\caption{Sensitivity of the HOMA-IR-by-sex interaction to exposure scale and upper-tail restriction}
\label{tab:s1_scale_tail}
\begin{threeparttable}
\small
\begin{tabular}{llrrrrr}
\toprule
HOMA scale & Analytic restriction & $n$ & Female slope & Male slope &
Male--female & Interaction $p$ \\
\midrule
Raw & Full sample & 3079 & 0.0629 & 0.0082 & $-0.0547$ & 0.00452 \\
Log & Full sample & 3079 & 0.2389 & 0.0166 & $-0.2223$ & 0.00205 \\
Raw & Below 99th percentile & 3048 & 0.0954 & 0.0219 & $-0.0735$ & 0.01334 \\
Log & Below 99th percentile & 3048 & 0.2440 & 0.0189 & $-0.2251$ & 0.00450 \\
Raw & Below 97.5th percentile & 3002 & 0.1188 & 0.0153 & $-0.1036$ & 0.00163 \\
Log & Below 97.5th percentile & 3002 & 0.2612 & 0.0003 & $-0.2609$ & 0.00135 \\
Raw & Below 95th percentile & 2925 & 0.1210 & 0.0209 & $-0.1001$ & 0.00946 \\
Log & Below 95th percentile & 2925 & 0.2432 & 0.0002 & $-0.2430$ & 0.00484 \\
\bottomrule
\end{tabular}
\begin{tablenotes}[flushleft]
\footnotesize
\item The log-HOMA coefficients are on a different exposure scale from the raw-HOMA
coefficients and should not be compared numerically as if they estimated the same
one-unit exposure contrast. All models retained $\ln(\mathrm{hsCRP})$ as the outcome
and the primary period-specific adjustment structure.
\end{tablenotes}
\end{threeparttable}
\end{table}

\begin{table}[htbp]
\centering
\caption{Sensitivity of the HOMA-IR-by-sex interaction to adiposity specification}
\label{tab:s2_adiposity}
\begin{threeparttable}
\small
\begin{tabular}{lrrrrr}
\toprule
Adiposity specification & $n$ & Female & Male &
Contrast & 95\% CI \\
\midrule
No adiposity adjustment & 3079 & 0.1892 & 0.0938 & $-0.0955$ & $-0.1616$ to $-0.0293$ \\
BMI & 3072 & 0.0606 & 0.0183 & $-0.0423$ & $-0.0811$ to $-0.0035$ \\
Waist circumference (primary) & 3079 & 0.0629 & 0.0082 & $-0.0547$ & $-0.0914$ to $-0.0180$ \\
Waist-to-height ratio & 3074 & 0.0522 & 0.0106 & $-0.0416$ & $-0.0791$ to $-0.0041$ \\
\bottomrule
\end{tabular}
\begin{tablenotes}[flushleft]
\footnotesize
\item Female and male columns report the estimated HOMA-IR slopes within
each sex; contrast is the male-minus-female slope difference.
Common-complete-case analyses among $n=3,072$ participants produced nearly
identical estimates, indicating that differences among adiposity specifications
were not attributable to differential missingness. These models are sensitivity
analyses of adjustment specification and are not mediation analyses.
\end{tablenotes}
\end{threeparttable}
\end{table}

\begin{landscape}
\begin{longtable}{p{2.0in}rrrrrr}
\caption{Targeted clinical adjustment and exclusion sensitivities}
\label{tab:s3_targeted}\\
\toprule
Model & $n$ & Female slope & Male slope & Male--female & $F$ & $p$ \\
\midrule
\endfirsthead
\toprule
Model & $n$ & Female slope & Male slope & Male--female & $F$ & $p$ \\
\midrule
\endhead
Smoking common-sample reference & 3075 & 0.0628 & 0.0081 & $-0.0548$ & 9.0765 & 0.00448 \\
+ smoking status & 3075 & 0.0632 & 0.0099 & $-0.0534$ & 8.6687 & 0.00537 \\
\addlinespace
Statin common-sample reference & 3041 & 0.0628 & 0.0080 & $-0.0548$ & 9.2612 & 0.00412 \\
+ statin use & 3041 & 0.0630 & 0.0081 & $-0.0548$ & 9.3704 & 0.00393 \\
\addlinespace
Renal common-sample reference & 3075 & 0.0625 & 0.0074 & $-0.0551$ & 8.9139 & 0.00481 \\
+ continuous eGFR & 3075 & 0.0623 & 0.0070 & $-0.0553$ & 9.0634 & 0.00450 \\
+ log creatinine & 3075 & 0.0623 & 0.0071 & $-0.0552$ & 9.0497 & 0.00453 \\
\addlinespace
ALT common-sample reference & 3074 & 0.0625 & 0.0075 & $-0.0551$ & 8.9101 & 0.00482 \\
+ log ALT & 3074 & 0.0633 & 0.0084 & $-0.0549$ & 8.8098 & 0.00504 \\
\addlinespace
Alcohol common-sample reference & 2919 & 0.0757 & 0.0084 & $-0.0673$ & 12.3816 & 0.00110 \\
+ alcohol-use category & 2919 & 0.0750 & 0.0085 & $-0.0665$ & 12.4028 & 0.00109 \\
\addlinespace
Exclude confirmed pregnancy & 3031 & 0.0666 & 0.0099 & $-0.0567$ & 9.0445 & 0.00454 \\
Exclude pregnancy + cannot determine & 3016 & 0.0668 & 0.0101 & $-0.0567$ & 8.9438 & 0.00475 \\
\bottomrule
\multicolumn{7}{p{8.4in}}{\footnotesize
Each added-covariate model should be compared with its reference model fitted
to the same participants. This is especially important for alcohol, for which
the common-sample reference contrast differs from the full-cohort primary
contrast because of sample restriction. eGFR used the 2021 CKD-EPI
creatinine equation. ALT and creatinine were natural-log transformed.}\\
\end{longtable}
\end{landscape}

\begin{table}[htbp]
\centering
\caption{Descriptive characterization by sex and hsCRP group}
\label{tab:s4_hscrp_descriptive}
\begin{threeparttable}
\small
\begin{tabular}{llrrrr}
\toprule
Measure & Statistic & Female $\leq10$ & Female $>10$ & Male $\leq10$ & Male $>10$ \\
\midrule
Unweighted $n$ & --- & 1510 & 128 & 1385 & 56 \\
HOMA-IR & Median & 1.91 & 3.01 & 2.05 & 2.04 \\
 & 75th percentile & 2.87 & 4.67 & 3.35 & 4.13 \\
 & 90th percentile & 4.51 & 7.40 & 5.13 & 6.77 \\
WBC ($10^3$ cells/$\mu$L) & Median & 6.2 & 7.5 & 6.3 & 7.2 \\
Absolute neutrophils ($10^3$ cells/$\mu$L) & Median & 3.5 & 4.45 & 3.5 & 4.5 \\
Neutrophils (\%) & Median & 57.4 & 60.0 & 56.6 & 61.4 \\
Absolute lymphocytes ($10^3$ cells/$\mu$L) & Median & 1.9 & 2.1 & 1.9 & 1.75 \\
Absolute monocytes ($10^3$ cells/$\mu$L) & Median & 0.5 & 0.5 & 0.5 & 0.7 \\
NLR & Median & 1.82 & 2.13 & 1.84 & 2.39 \\
\bottomrule
\end{tabular}
\begin{tablenotes}[flushleft]
\footnotesize
\item hsCRP is expressed in mg/L. CBC measures were available for all 184
participants with hsCRP $>10$ mg/L. Values are descriptive and were not used
to diagnose acute infection, assign an inflammatory etiology, or define
additional exclusions. WBC = white blood cell count; NLR =
neutrophil-to-lymphocyte ratio.
\end{tablenotes}
\end{threeparttable}
\end{table}


\end{document}
